\begin{frame}{Thuật toán DFS (Depth-First Search)}
    \begin{itemize}
        \item \textbf{Nguyên lý:} "Đi sâu nhất có thể" trước khi quay lui (Backtracking).
        \item \textbf{Cấu trúc dữ liệu:} Sử dụng \textbf{Ngăn xếp (Stack)} hoặc \textbf{Đệ quy} (chính là Stack của hệ thống).
        \item \textbf{Các bước:}
        \begin{enumerate}
            \item Thăm đỉnh $u$, đánh dấu $u$ là đã thăm (visited).
            \item Với mỗi đỉnh $v$ kề với $u$:
            \begin{itemize}
                \item Nếu $v$ chưa được thăm, thực hiện gọi đệ quy \texttt{DFS(v)}.
            \end{itemize}
        \end{enumerate}
    \end{itemize}
    \vspace{0.2cm}
    \textit{Lưu ý:} DFS rất dễ cài đặt nhờ đệ quy, nhưng cần cẩn thận với đồ thị quá lớn vì có thể gây tràn Stack (Stack Overflow).
\end{frame}
\begin{frame}[fragile]{Cài đặt DFS với C++ (Đệ quy)}
\begin{verbatim}
vector<int> adj[100]; // Danh sách kề
bool visited[100];    // Mảng đánh dấu đã thăm

void DFS(int u) {
    visited[u] = true;      // Đánh dấu u đã thăm
    // Duyệt qua các đỉnh kề v
    for(int v : adj[u]) {
        if(!visited[v]) {
            DFS(v);         // Gọi đệ quy cho v
        }
    }
}
\end{verbatim}
\textbf{Ghi chú:} visited[] đóng vai trò cực kỳ quan trọng để ngăn chương trình lặp vô tận (nếu đồ thị có chu trình).
\end{frame}